% chktex-file 8
% chktex-file 12
% chktex-file 18
\beginsong{Pražce}[by={Pavel Dobeš}]

\beginverse{}
\ch{C}{}{}{}Házím tornu na svý záda, feldflašku a \ch{G7}{}{}{}sumky,
navštívím dnes kamaráda z železniční průmky.
\endverse{}

\beginchorus{}
Vždyť je j\ch{C}{}{}{}aro, zapni si kšandy,
pozdravuj \ch{G7}{}{}{}vlaštovky a, muziko, ty \ch{C}{}{}{}hraj.
\endchorus{}

\beginverse{}
Vystupuji z vlaku, který mizí v dálce,
stojím v České Třebové a všude kolem pražce.
\endverse{}

\beginverse{}
Pohostil mě slivovicí, představil mě Mařce,
posadil mě na lavici z dubového pražce.
\endverse{}

\beginverse{}
Provedl mě domem -- nikde kousek zdiva,
všude samej pražec, jen Máňa byla živá.
\endverse{}

\beginchorus{}
To je to jaro, zapni si kšandy,
pozdravuj vlaštovky a, muziko, ty hraj.
\endchorus{}

\beginverse{}
Plakáty nás informují:"Přijď pracovat k dráze,
pakliže ti vyhovují rychlost, šmír a saze."
\endverse{}

\beginverse{}
A jestliže jsi labužník a přes kapsu se praštíš,
upečeš i krávu na železničních pražcích.
\endverse{}

\beginverse{}
A naučíš se skákat tak, jak to umí vrabec,
když na nohu si pustíš železniční pražec.
\endverse{}

\beginverse{}
Když má děvče z Třebové rádo svého chlapce,
posílá mu na vojnu železniční pražce.
\endverse{}

\beginverse{}
A když děti zlobí, tak hned je doma mazec,
Děda Mráz jim nepřinese ani jeden pražec.
\endverse{}

\beginverse{}
Před děvčaty z Třebové chlubil jsem se silou,
pozvedl jsem pražec, načež odvezli mě s kýlou.
\endverse{}

\beginverse{}
Pamatuji pouze ještě operační sál,
pak praštili mě pražcem a já jsem tvrdě spal.
\endverse{}

\beginchorus{}
A bylo jaro, zapni si kšandy,
lítaly vlaštovky a zelenal se háj.
\endchorus{}

\endsong{}
